% Gerado a partir de docs/tcc.md por scripts/sincronizar_tcc.py.
\chapter{Metodologia}\label{metodologia}

Este capítulo apresenta o caminho seguido para transformar a ideia do sistema de avaliação contextual em um protótipo demonstrável. A metodologia combinou pesquisa aplicada, revisão bibliográfica e normativa, Design Science Research, prototipação incremental e Lean Inception, de modo que a definição do produto e a implementação técnica avançassem de forma coerente.

\section{Caracterização da Pesquisa}\label{caracterizauxe7uxe3o-da-pesquisa}

A pesquisa é aplicada e exploratória: propõe uma solução computacional para relacionar contexto de navegação e avaliação declarada em serviços públicos digitais. O artefato possui dois componentes: uma implementação funcional local em Python, com API, páginas, formulário e painel; e um protótipo de interface no Figma, com os fluxos A e B usados no estudo com participantes. A implementação recebeu inspeção interna funcional e heurística. As interfaces no Figma receberam avaliação exploratória de usabilidade, instrumentada pelo Useberry. Nenhum desses procedimentos constitui validação em contexto real de prestação de serviço público.

O percurso dialoga com a Design Science Research, na qual a construção e a avaliação de artefatos apoiam a produção de conhecimento em sistemas de informação \cite{hevner2004designscience}. A demonstração técnica da implementação e as observações de interface oferecem evidências de naturezas diferentes. A avaliação com oito participantes se refere às interfaces no Figma; não verifica os endpoints, a persistência ou os gatilhos da implementação Python. A estruturação completa do percurso pelas etapas da DSR permanece parte da continuidade metodológica.

As fontes foram diferenciadas por natureza e finalidade. Artigos e livros acadêmicos fundamentam conceitos e questões de pesquisa; normas oferecem definições e critérios; legislação caracteriza o enquadramento normativo; relatórios técnicos e materiais institucionais descrevem modelos e práticas de governo digital. Os módulos da Enap são material didático institucional. Caroli, Fowler e Caroli, Cooper e Kalbach são usados como literatura profissional de concepção e projeto; Reichheld caracteriza a proposta original de NPS, sem constituir validação desse indicador para serviços públicos. Hotjar, Microsoft Clarity e Useberry são documentação comercial, limitada à descrição de funcionalidades declaradas e instrumentos utilizados. A classificação das fontes é apresentada a seguir; a presença de uma fonte não comprova a eficácia do mecanismo proposto.

\begin{itemize}
\item
  \textbf{Legislação:} \citeonline{brasil2011lei12527}; \citeonline{brasil2017decreto9094}; \citeonline{brasil2017lei13460}; \citeonline{brasil2018lgpd}; \citeonline{brasil2021lei14129}; \citeonline{brasil2022decreto11260}; \citeonline{brasil2022portaria548}.
\item
  \textbf{Documentação ou pesquisa institucional:} \citeonline{govbrFerramentaAvaliacao}; \citeonline{brasilMgi2024ppaParticipativo}; \citeonline{cetic2025ticDomicilios2024}; \citeonline{govuk2021completionRate}; \citeonline{govuk2021usingData}.
\item
  \textbf{Literatura profissional:} \citeonline{caroli2018leanInception}; \citeonline{cooper1999inmates}; \citeonline{fowler2022leanInception}; \citeonline{kalbach2016mapping}; \citeonline{reichheld2003nps}.
\item
  \textbf{Literatura acadêmica/científica:} \citeonline{deloneMclean2003}; \citeonline{gronroos1984service}; \citeonline{hektner2007experienceSampling}; \citeonline{hevner2004designscience}; \citeonline{kahneman1993peakEnd}; \citeonline{larson1983experienceSampling}; \citeonline{manfreda2008webSurveys}; \citeonline{mark2008interruptedWork}; \citeonline{nielsen1994heuristic}; \citeonline{oliver2014satisfaction}; \citeonline{parasuraman1988servqual}; \citeonline{schwarz1999selfReports}; \citeonline{shihFan2008responseRates}; \citeonline{Sommerville2011}; \citeonline{sweller1988cognitiveLoad}.
\item
  \textbf{Documentação comercial:} \citeonline{hotjar2026}; \citeonline{microsoftClarity2026}; \citeonline{useberryFigma}.
\item
  \textbf{Norma técnica:} \citeonline{iso9241112018}; \citeonline{iso92412102019}.
\item
  \textbf{Material didático institucional:} \citeonline{limaPereira2020modulo1}; \citeonline{limaPereira2020modulo2}; \citeonline{limaPereira2020modulo3}; \citeonline{limaPereira2020modulo4}.
\item
  \textbf{Relatório técnico institucional:} \citeonline{pedrosa2019qualidadeGovbr}; \citeonline{silva2024atributosApi}; \citeonline{silva2019brasp}; \citeonline{silva2022modeloUnificado}.
\end{itemize}

\section{Lean Inception}\label{lean-inception}

A Lean Inception foi utilizada como protocolo metodológico para alinhar visão de produto, escopo e MVP. A abordagem de \citeonline{caroli2018leanInception} auxilia equipes a definir o produto certo antes de ampliar o desenvolvimento; \citeonline{fowler2022leanInception} também a descreve como uma inception focada em compreender funcionalidades, clientes e Canvas MVP. Neste trabalho, a técnica organizou a definição do sistema de avaliação contextual e ajudou a transformar a ideia inicial em um conjunto priorizado de funcionalidades.

A atividade interna contou com quatro participantes. As perspectivas listadas a seguir representam suas contribuições à discussão, sem atribuir-lhes vínculo institucional ou condição de especialista independente:

\begin{itemize}
\item
  \textbf{Felipe}, atendimento público, representando dúvidas recorrentes e dificuldades relatadas por cidadãos;
\item
  \textbf{Lucas}, gestão digital, representando a necessidade de indicadores acionáveis para priorização de melhorias;
\item
  \textbf{Diassis}, estudante da UnB, representando viabilidade técnica, escopo acadêmico, segurança e LGPD;
\item
  \textbf{Bruno}, estudante da UnB, representando experiência do usuário, jornada, clareza de interface e validação.
\end{itemize}

Lucas, participante da discussão, e Lucas Almeida, persona fictícia da seção 4.3, são referências distintas. A participação de Wesley Pedrosa dos Santos na Lean Inception esteve voltada ao tema de seu próprio trabalho; não é apresentada como avaliação independente do sistema aqui proposto. A dinâmica foi uma atividade interna de alinhamento e delimitação do produto, não uma validação externa. As decisões da Lean Inception não substituem especificação de requisitos, teste de usabilidade ou avaliação por especialistas independentes.

\section{Etapas da Lean Inception}\label{etapas-da-lean-inception}

A agenda da Lean Inception foi organizada para que as funcionalidades surgissem a partir da jornada do usuário, e não apenas de possibilidades técnicas. As etapas executadas foram:

\begin{enumerate}
\def\labelenumi{\arabic{enumi}.}
\item
  \textbf{Kick-off}: alinhamento do problema, da hipótese e dos limites do MVP;
\item
  \textbf{É / Não É / Faz / Não Faz}: delimitação do produto como solução de feedback contextual, não como ouvidoria, chatbot, ferramenta de identificação individual ou sistema governamental completo;
\item
  \textbf{Objetivos do Produto}: definição de objetivos como capturar sinais de possível dificuldade, solicitar feedback no momento adequado, preservar privacidade e gerar dados acionáveis;
\item
  \textbf{Personas}: definição de Carlos Pereira, cidadão usuário de serviço público digital, e Lucas Almeida, gestor público de serviço digital;
\item
  \textbf{Jornadas}: comparação entre um fluxo com formulário estático ao final e um fluxo com feedback contextual durante a dificuldade;
\item
  \textbf{Brainstorming de features}: levantamento de funcionalidades relacionadas a coleta de sinais, regras, feedback, painel e integração;
\item
  \textbf{Revisão técnica, de negócio e UX}: avaliação relativa de esforço, valor para gestão e impacto na experiência do usuário;
\item
  \textbf{Sequenciador}: organização das funcionalidades em ondas de entrega;
\item
  \textbf{Canvas MVP}: consolidação da menor versão viável para validação inicial.
\end{enumerate}

A construção de personas foi baseada na tradição de design centrado no usuário. \citeonline{cooper1999inmates} apresenta personas como arquétipos usados para orientar decisões de produto. O mapeamento de jornada seguiu a lógica de visualização de experiências discutida por \citeonline{kalbach2016mapping}, com identificação de pontos de contato, fricções e oportunidades de intervenção. As personas também consideraram o cenário brasileiro de uso de tecnologias digitais, descrito pela TIC Domicílios 2024 \cite{cetic2025ticDomicilios2024}, e a existência de mecanismos oficiais de avaliação de serviços públicos digitais \cite{govbrFerramentaAvaliacao}.

\section{Metodologia do Desenvolvimento da Solução}\label{metodologia-do-desenvolvimento-da-soluuxe7uxe3o}

O desenvolvimento adotou prototipação incremental, com a finalidade de explorar decisões de solução e reduzir incertezas \cite{Sommerville2011}. A Lean Inception foi conduzida entre 18 e 27 de junho de 2026; a avaliação das interfaces no Figma ocorreu em 28 de junho de 2026. A existência de uma implementação funcional em Python antes das sessões não significa que ela tenha sido usada pelos participantes. Não se atribuem os resultados do Figma à API.

O sequenciador de concepção organizou cinco ondas:

\begin{enumerate}
\def\labelenumi{\arabic{enumi}.}
\tightlist
\item
  \textbf{Onda 1 --- Base técnica e desenho comparativo:} identificador de sessão, eventos e planejamento dos formulários contextual e final.
\item
  \textbf{Onda 2 --- Regras e feedback:} condições de acionamento, sugestão de redirecionamento e formulário.
\item
  \textbf{Onda 3 --- Painel básico:} agregações de eventos e respostas e planejamento de indicadores comparativos.
\item
  \textbf{Onda 4 --- Incrementos de gestão:} parâmetros por serviço, filtros e apresentação de comentários.
\item
  \textbf{Onda 5 --- Preparação de integração real:} organizações, autenticação, chave de API e cliente instalável.
\end{enumerate}

As ondas 1 a 3 constituíram o MVP planejado, mas foram implementadas apenas em parte. Na aplicação Python existem identificadores técnicos, captura de cliques em elementos instrumentados, envio de duração simulada, regras, formulário reutilizável e agregações básicas. Não foram implementados cronômetro automático de navegação, registro dos grupos A/B, exposição ao convite, ranking por alertas ou taxa de resposta por fluxo. A comparação de interfaces foi realizada no Figma. Parâmetros por serviço, originalmente previstos como incremento da onda 4, já estão no backend; filtros e comentários recentes continuam pendentes. O sequenciador expressa o planejamento, não uma lista de entregas integralmente executadas.

\section{Hipótese e Critérios da Avaliação Preliminar}\label{hipuxf3tese-e-crituxe9rios-da-avaliauxe7uxe3o-preliminar}

As hipóteses direcionais propostas para um estudo futuro são: \textbf{H1}, em tarefas e condições equivalentes, o fluxo contextual terá maior proporção de participantes que enviam feedback do que o fluxo com formulário final; \textbf{H2}, os comentários enviados no fluxo contextual terão maior escore de especificidade do que os enviados no fluxo com formulário final, segundo rubrica definida antes da coleta. As hipóteses não foram testadas por esta avaliação exploratória. A regra de exposição ao convite e os denominadores deverão ser definidos previamente, com divulgação das respostas tanto por participante alocado quanto por participante exposto.

As variáveis previstas têm diferentes condições de observação nesta etapa:

\begin{itemize}
\tightlist
\item
  \textbf{Realização da avaliação:} considerou-se resposta, no fluxo A, quando o participante realizou a avaliação apresentada ao final da jornada e, no fluxo B, quando aceitou o convite e realizou a avaliação. Abrir o convite, sem realizar a avaliação, não corresponde a esse critério. Essa classificação não comprova envio de dados à API nem digitação de comentário. A exposição ao formulário final não foi discriminada para cada participante do A; as proporções por participante e por exposição devem permanecer distintas.
\item
  \textbf{Conclusão da tarefa:} há classificação Sim/Parcial/Não; o resumo considera somente Sim como conclusão completa. Isso explicita o cálculo descritivo e não representa critério prévio de sucesso documentado.
\item
  \textbf{Esforço percebido:} não foi medido por instrumento específico; a nota de facilidade do Apêndice F não corresponde ao CES.
\item
  \textbf{Qualidade/especificidade dos comentários do formulário:} não há corpus separado nem rubrica aplicada. Os comentários do Apêndice G são do questionário pós-sessão e não permitem testar H2.
\item
  \textbf{Interrupção e abandono associados ao convite:} há relatos pontuais de ignorar o convite; não há medida capaz de atribuir causalidade ao popup ou ao feedback.
\item
  \textbf{Adequação do momento:} existem relatos pontuais, mas o conjunto de dados apresentado não inclui a série completa de respostas à pergunta correspondente do questionário.
\end{itemize}

H1 e H2 são as hipóteses comparativas principais. Conclusão da tarefa, esforço, interrupção e adequação do momento são variáveis secundárias previstas para caracterizar benefícios e custos da intervenção; nesta etapa, sua observação foi incompleta. H2 compara somente comentários de quem respondeu ao formulário. Essa seleção pode diferir entre os grupos e produzir viés; um estudo futuro deverá relatar a ausência de respostas e a composição dos respondentes, sem generalizar os escores para todos os participantes.

Tempo e cliques são medidas auxiliares. A tarefa no Figma foi instrumentada pelo Useberry, enquanto as regras da API constituem outra fonte de evidência, relativa à implementação. O planejamento de instrumentos e critérios comparativos deve preceder nova coleta.

Para a inspeção interna do artefato foram organizados sete critérios, retomados pelos mesmos identificadores na Tabela \ref{tab:criterios-avaliacao-preliminar}:

\begin{enumerate}
\def\labelenumi{\arabic{enumi}.}
\tightlist
\item
  \textbf{C1 --- Registro mínimo de navegação:} existência de sessão técnica e estruturas para registrar eventos, com explicitação dos limites de captura do cliente de demonstração.
\item
  \textbf{C2 --- Acionamento contextual:} correspondência entre parâmetros por serviço e ações previstas no código, sem interpretar acionamento como dificuldade confirmada.
\item
  \textbf{C3 --- Esforço de resposta:} presença de formulário curto como decisão de projeto; esforço efetivamente percebido deve ser tratado como não medido nesta coleta.
\item
  \textbf{C4 --- Apresentação para gestão:} existência de agregações e alertas; utilidade e acionabilidade dependem de avaliação com gestores.
\item
  \textbf{C5 --- Inspeção das regras:} regras localizáveis no código e parâmetros por serviço; histórico de configuração e versionamento permanecem ausentes.
\item
  \textbf{C6 --- Privacidade:} exame dos campos e das limitações de transparência, retenção e acesso, sem concluir conformidade jurídica.
\item
  \textbf{C7 --- Interrupção:} possibilidade de ignorar o convite e limitar o popup; efeito sobre abandono ou irritação depende de mensuração posterior.
\end{enumerate}

A apresentação dos resultados como exploratórios é uma regra de interpretação comum a todos os critérios, sem afirmação de superioridade ou generalização.

\section{Procedimentos da Avaliação Exploratória}\label{procedimentos-da-avaliauxe7uxe3o-exploratuxf3ria}

O relato do trabalho distingue três procedimentos. O primeiro é a inspeção funcional interna da implementação, com observação do registro de eventos, regras e indicadores descritos pelos autores, sem relatório de execução de testes anexado. O segundo é uma \textbf{avaliação heurística interna/preliminar}, também dos autores. O terceiro é a avaliação de fluxos de interface no Figma com oito participantes; os resultados desse procedimento não devem ser apresentados como verificação da API.

A avaliação exploratória de interface ocorreu em 28 de junho de 2026, com oito participantes externos ao grupo de desenvolvimento, quatro em cada fluxo. A alocação entre os grupos foi auxiliada por inteligência artificial e descrita como sorteio, sem registro auditável do mecanismo. Não se assume aleatorização verificável. O Figma foi instrumentado pelo Useberry para registro automático de cliques e tempos. O fluxo A previa formulário final e o B convite contextual; o popup de redirecionamento fez parte de ambos. A tarefa consistiu em localizar e iniciar a declaração do Imposto sobre a Renda da Pessoa Física (IRPF) de 2026, com limite de dez minutos. As sessões foram individuais e moderadas, sob coordenação de Julio Cesar Gomes de Souza Crispim e Fernando Gabriel dos Santos Carrijo, com TCLE que mencionava o Useberry e o registro de interações. A gravação se limitou à tela e aos eventos de interação, sem áudio ou vídeo de câmera dos participantes. Os comentários apresentados no Apêndice G provêm do questionário pós-sessão. O convite no fluxo B podia decorrer de tempo, entrada em tela ou clique no Figma. No fluxo A, resposta significa realizar a avaliação no formulário final; no B, aceitar o convite e realizar a avaliação sobre a experiência, não apenas abrir o convite. Os valores temporais, as telas e as interações específicas não estão discriminados neste protocolo. Esse procedimento permite observação formativa de interface, sem estabelecer funcionamento dos gatilhos do backend ou efeitos causais entre fluxos.

\section{Tratamento dos Dados}\label{tratamento-dos-dados}

A implementação organiza dois grupos de dados: eventos de interação, com contexto de sessão, página, serviço, clique e duração informada; e respostas declaradas, com estrelas, atributos, pergunta auxiliar de recomendação e comentário opcional. Eventos indicam contexto e possível esforço, sem medir qualidade percebida. Estrelas e atributos expressam julgamentos declarados; a pergunta NPS não foi validada como medida de qualidade percebida neste cenário. A base do estudo de interface é distinta: cliques e tempos registrados no Useberry, classificações de tarefa e respostas ao questionário pós-sessão, transcritos no Apêndice G. Não foi apresentado um conjunto separado de respostas do formulário que permita comparar seus comentários, estrelas e atributos.

A implementação oferece consulta de alertas, cliques e identificadores de sessão por serviço, total de feedbacks e médias de respostas. Não oferece contagem de sessões únicas alertadas, recorrência de atributos ou análise A/B. O formulário não solicita nome, CPF, e-mail ou credencial de login. O backend aceita um \texttt{user\_id} opcional e registra URL e coordenadas do clique na tela, não geolocalização. Essas condições não comprovam anonimização. Aviso de coleta, retenção e controle de acesso exigem requisitos próprios \cite{brasil2018lgpd}.
